% Sección aditiva. No modifica la edición de 211 páginas.
% Procedencia y comprobaciones: documentos internos de este mismo directorio.
\section{Refinamiento nonádico, banderas positivas y compatibilidad de los momentos}
\label{sec:dim-cofinal}

La longitud de una publicación finita determina cuántos menores pueden
formarse con sus columnas. El refinamiento conserva las columnas heredadas
y añade otras con procedencia conocida. Por ello, la dimensión de una
realización matricial puede crecer sin sustituir el registro que fija sus
separaciones. La construcción siguiente reúne este crecimiento con dos
nociones precisas de compatibilidad: restricción de evaluaciones a las
marcas anteriores y promedio de observables sobre las celdas descendientes.
La primera preserva las banderas polinomiales; la segunda conserva una
medida y separa el incremento de información en detalles ortogonales.

El dato de partida es la salida dodecafásica ya producida por
APP--TRIT--TPK. Las dos hojas de APP conservan sus divisiones completas;
TRIT orienta el régimen y el operador
\(U_t=\operatorname{Upd}_t\circ\operatorname{Tra}_t\circ
\operatorname{Sel}_t\) transporta el estado enriquecido. La prolongación
\(w_6\to w_{12}\to w_{18}\to w_{24}\to w_{30}\to R_{36}\to G_9\)
conserva retorno de fase y avance de memoria. La estructura discreta conjunta
del continuo conserva correlativamente las publicaciones de ese estado,
con sus hojas, orientación, memoria y fronteras. De ese mismo estado proceden
las publicaciones regionales y el registro firmado cuya inversión da
\[
 K=(234,543,140,729,659,824,621,58,914,794,146,601),
 \qquad Q=\sum_{j=1}^{12}K_j=6263.
\]
Se utiliza esta salida con sus marcas y orientación; las matrices de
momentos son realizaciones posteriores. El índice de refinamiento que se
introduce abajo cuenta subdivisiones del lector geométrico. Su identificación
con una duración física o con un número de microtransiciones del TPK
requiere conservar el reloj de esa otra realización.

\subsection{Particiones marcadas y crecimiento cofinal}

Escribamos \(d_j=K_j/Q\), \(a_0=0\) y
\(a_j=\sum_{h=1}^j d_h\). Para \(r\geq0\), definimos
\[
 L_r=12\,9^r,\qquad N_r=L_r+1,\qquad
 t_{r,(j-1)9^r+c}=a_{j-1}+\frac{c}{9^r}d_j,
 \quad 0\leq c\leq9^r.
\]
Los extremos comunes se cuentan una sola vez. Las \(L_r\) celdas se
indexan por \((j,c)\), con \(1\leq j\leq12\) y
\(0\leq c<9^r\), y tienen longitud
\(\ell_{r,j,c}=d_j/9^r\). Se conserva también la escritura de \(c\)
en \(r\) cifras nonádicas, incluidos los ceros iniciales. Un hijo tiene
índice \(9c+h\), con \(0\leq h\leq8\); la división euclídea recupera
su padre y la última cifra. Los extremos compartidos retienen su marca
de frontera; la convención de intervalos semiabiertos sólo organiza
las celdas y no elimina ninguna marca.

\begin{proposition}[Cofinalidad y recuperación del registro]
\label{prop:dim-grid}
Las marcas satisfacen
\[
 0=t_{r,0}<t_{r,1}<\cdots<t_{r,L_r}=1,
 \qquad t_{r+1,9i}=t_{r,i}.
\]
El máximo de las longitudes es
\(h_r=(\max_jK_j)/(Q9^r)\), que tiende a cero. Para todo par de
enteros \(k\geq1\), \(m\geq1\), existe \(r_0\) tal que
\(k+m\leq N_r\) para \(r\geq r_0\). El registro se recupera en
cualquier nivel por
\[
 K_j=Q\sum_{c=0}^{9^r-1}\ell_{r,j,c}.
\]
\end{proposition}
\begin{proof}
Cada diferencia consecutiva es \(d_j/9^r>0\). Al reemplazar \(c\)
por \(9c\), la fórmula del nivel siguiente da exactamente la marca
anterior. Las nueve longitudes de los hijos suman la longitud del padre;
la iteración prueba la recuperación. Finalmente \(N_r=12\,9^r+1\)
crece sin cota. Una elección explícita es cualquier \(r_0\geq0\) con
\(9^{r_0}\geq(k+m-1)/12\). Todo el argumento conserva \(Q\), el
índice de bloque y las cifras de descendencia.
\end{proof}

La normalización determina las razones \(K_j/Q\). La recuperación de
los enteros utiliza, además, la carga \(Q\) conservada por el registro.
Así se distinguen la geometría normalizada y la información de escala
que acompaña a su publicación.

\subsection{Banderas de evaluaciones y realizaciones en toda dimensión finita}

Para \(1\leq q\leq N_r\), sea
\[
 V_r^{(q)}=(t_{r,i}^{\,a})_{0\leq a<q,\,0\leq i<N_r},
 \qquad F_r^q=\operatorname{fila}V_r^{(q)}.
\]
La potencia de exponente cero vale uno también en el nodo cero.

\begin{theorem}[Bandera positiva y composición de rango arbitrario]
\label{thm:dim-all-ranks}
Los espacios \(F_r^q\) forman una bandera
\(F_r^1\subset F_r^2\subset\cdots\subset F_r^{N_r}\), con
\(\dim F_r^q=q\), y cada matriz \(V_r^{(q)}\) tiene todos sus
menores máximos ordenados estrictamente positivos. Para
\(k\geq1\), \(m\geq1\), \(k+m\leq N_r\), las matrices
\[
 C_r=V_r^{(k)},\quad Z_r=V_r^{(k+m)},\quad
 Y_r=C_rZ_r^{\mathsf T}
\]
satisfacen
\[
 (Y_r)_{ab}=\sum_{i=0}^{L_r}t_{r,i}^{a+b},\qquad
 \operatorname{rank}Y_r=k,\qquad
 [Y_r]\in\mathcal A_{N_r,k,m}(Z_r).
\]
La construcción alcanza todo par finito \((k,m)\) mediante un nivel
suficientemente profundo de la misma partición marcada.
\end{theorem}
\begin{proof}
Para \(i_1<\cdots<i_q\), el menor correspondiente es
\[
 \det(t_{r,i_b}^{a-1})_{a,b=1}^{q}
   =\prod_{a<b}(t_{r,i_b}-t_{r,i_a})>0.
\]
Esto prueba rango, positividad e inclusiones de la bandera. El bloque
inicial \(k\times k\) de \(Y_r\) es el Gram de las evaluaciones
de \(1,t,\ldots,t^{k-1}\). Para \(x\neq0\),
\[
 x^{\mathsf T}(Y_r)_{[0,k-1],[0,k-1]}x
 =\sum_i\left(\sum_{a=0}^{k-1}x_at_{r,i}^{a}\right)^2>0,
\]
porque un polinomio no nulo de grado inferior a \(k\) no se anula
en los \(N_r\geq k+m\) nodos distintos. El bloque es invertible,
y \(Y_r\), que tiene \(k\) filas, tiene rango \(k\).
La pertenencia indicada es la aplicación de la matriz externa positiva
\(Z_r\) al punto positivo representado por \(C_r\).
\end{proof}

Se emplea la convención matemática
\[
 \mathcal A_{N,k,m}(Z)=
 \{[CZ^{\mathsf T}]:[C]\in\operatorname{Gr}_{\geq0}(k,N)\}.
\]
La definición para \(m\) general puede consultarse en S.~N.~Karp,
L.~K.~Williams y Y.~X.~Zhang,
\href{https://people.math.harvard.edu/~williams/papers/amplituhedron-plus.pdf}
{\emph{Decompositions of amplituhedra}}, definición~1.1.
Aquí el dato externo ya procede de la partición publicada por HMT.
La conclusión construye una familia determinada de puntos y banderas;
las dimensiones del Grassmanniano ambiente y de un atlas de su imagen
son objetos adicionales a los rangos probados.

Sea \(E_r:\mathbb R^{N_{r+1}}\to\mathbb R^{N_r}\) la restricción
a las posiciones \(9i\). Para todo \(q\leq N_r\),
\[
 E_r\bigl((V_{r+1}^{(q)})^{\mathsf T}x\bigr)
       =(V_r^{(q)})^{\mathsf T}x.
\]
Por tanto, \(E_r:F_{r+1}^q\to F_r^q\) es un isomorfismo sobre
estos espacios de evaluaciones: sus dos realizaciones identifican el
mismo polinomio de grado inferior a \(q\). La composición de las
restricciones coincide con la restricción directa. Los menores sobre
índices heredados conservan exactamente su valor.

\begin{proposition}[Recuperación desde una bandera calibrada]
Si se conservan el orden de las columnas, las marcas extremas y la
calibración \(t_{r,0}=0\), \(t_{r,L_r}=1\), el plano \(F_r^2\)
determina los nodos y las separaciones. Junto con \(Q\) y las
marcas de descendencia determina \(K\).
\end{proposition}
\begin{proof}
El plano contiene el vector constante y el vector de nodos. Cualquier
otra coordenada que, junto con el vector constante, genere el mismo
plano es \(a+bt\), con \(b\neq0\). Las dos calibraciones fijan
\(a=0\), \(b=1\). Las diferencias consecutivas recuperan las
longitudes, y la proposición~\ref{prop:dim-grid} recupera el registro.
\end{proof}

\subsection{Cuerpos convexos de rango uno y crecimiento del dominio}

Para \(m\geq1\) y \(m+1\leq N_r\), definimos
\[
 P_r^{(m)}=\operatorname{conv}
 \{(t_{r,i},t_{r,i}^2,\ldots,t_{r,i}^m):0\leq i\leq L_r\}.
\]
La normalización de una fila positiva de \(C\) como coeficientes
baricéntricos prueba
\(\mathcal A_{N_r,1,m}(Z_r)=P_r^{(m)}\). El determinante de
Vandermonde proporciona \(m+1\) puntos afínmente independientes,
de modo que el cuerpo tiene dimensión \(m\).

\begin{proposition}[Cobertura simplicial a dimensión finita y límite convexo]
\label{prop:dim-polytopes}
Cada \(P_r^{(m)}\) admite una triangulación por sus vértices, con
grado uno sobre el interior de cada celda simplicial. La suma de sus
formas canónicas tiene residuos interiores cancelados. Además,
\[
 P_r^{(m)}\subseteq P_{r+1}^{(m)},\qquad
 \overline{\bigcup_{r\geq r_0}P_r^{(m)}}=
 \operatorname{conv}\{(t,t^2,\ldots,t^m):0\leq t\leq1\}.
\]
\end{proposition}
\begin{proof}
Se ordenan los vértices por sus marcas. En cada cara se elige su menor
vértice; se triangulan recursivamente las facetas que no lo contienen
y se forman los conos correspondientes. La inducción sobre la dimensión
prueba cobertura e interiores disjuntos, y la misma regla en caras
comunes garantiza compatibilidad. En cada simplex orientado
\([v_0,\ldots,v_m]\), con matriz
\(A=(v_1-v_0,\ldots,v_m-v_0)\) de determinante positivo y
coordenadas baricéntricas \(\lambda_0,\ldots,\lambda_m\), la
aplicación baricéntrica es afín e invertible. Su forma es
\[
 \Omega_{[v_0,\ldots,v_m]}=
 \frac{dx_1\wedge\cdots\wedge dx_m}
 {\det A\prod_{j=0}^{m}\lambda_j}.
\]
La restricción residual a cada cara es su forma orientada. Dos simplices
adyacentes inducen orientaciones opuestas sobre su cara común, de modo
que los residuos interiores se cancelan. Esta es la misma operación de
triangulación y residuo aplicada en la realización de rango uno anterior,
con el número de vértices ahora aumentado por el refinamiento.

La inclusión de nodos heredados prueba la inclusión de cuerpos. La malla
tiende a cero, por lo que las curvas muestreadas son densas en la curva
compacta de momentos. Toda combinación convexa finita de puntos de la
curva se aproxima por combinaciones con nodos de un nivel común.
Recíprocamente, todos los cuerpos están contenidos en la envolvente
convexa compacta de la curva. Ambas inclusiones prueban la igualdad.
\end{proof}

La cancelación de residuos opera sobre una triangulación o subdivisión
de un cuerpo fijo. El paso de \(P_r^{(m)}\) a \(P_{r+1}^{(m)}\)
añade vértices y puede ampliar el cuerpo. Sus formas pertenecen a esos
dominios respectivos. El límite convexo precedente especifica la
convergencia de los cuerpos, sin imponer una identidad entre sus formas.

\subsection{Medida del calendario y convergencia de los momentos nodales}

Sea \(F_K:[0,1]\to[0,1]\) la función creciente afín por tramos
que lleva \((j-1)/12\) a \(a_{j-1}\) y \(j/12\) a \(a_j\).
En el intervalo correspondiente su pendiente es \(12d_j>0\).
La medida del calendario transportado es
\[
 \mu_K=(F_K)_*(ds),\qquad
 \frac{d\mu_K}{dt}=\frac1{12d_j}\quad
 (a_{j-1}<t<a_j).
\]
Cada intervalo inicial tiene masa \(1/12\), y cada celda de nivel
\(r\) tiene masa \(1/L_r\). Esta elección asigna el mismo peso
a las posiciones del calendario antes de realizar sus longitudes.
La medida de longitud \(dt\) es otro lector; ambas coinciden cuando
todas las separaciones iniciales son iguales.

Sus momentos son racionales y están determinados por \(K\):
\begin{equation}
 M_p=\int_0^1t^p\,d\mu_K(t)
   =\frac1{12(p+1)}\sum_{j=1}^{12}
       \frac{a_j^{p+1}-a_{j-1}^{p+1}}{d_j},\qquad p\geq0.
 \label{eq:dim-moments}
\end{equation}
\begin{proposition}[Recuperación por momentos y cuantiles]
\label{prop:dim-moment-recovery}
La sucesión completa \((M_p)_{p\geq0}\) determina \(\mu_K\).
Si \(H_K(t)=\mu_K([0,t])\) es su función de distribución, entonces
\[
 a_j=H_K^{-1}(j/12),\qquad
 K_j=Q\bigl(H_K^{-1}(j/12)-H_K^{-1}((j-1)/12)\bigr).
\]
En consecuencia, los momentos completos, la carga \(Q\) y el orden
del calendario recuperan el registro inicial.
\end{proposition}
\begin{proof}
Dos probabilidades sobre \([0,1]\) con los mismos momentos tienen las
mismas integrales para todos los polinomios. La aproximación uniforme
de funciones continuas por polinomios y la masa finita extienden esa
igualdad a todas las funciones continuas; éstas determinan la medida
de Borel. La densidad de \(\mu_K\) es positiva y constante en cada
intervalo inicial. Su distribución es, por tanto, continua y
estrictamente creciente, y cumple \(H_K(a_j)=j/12\). La inversión
de esta igualdad y la diferencia de extremos prueban las fórmulas.
\end{proof}

La probabilidad nodal correspondiente al teorema
\ref{thm:dim-all-ranks} es
\(\nu_r=N_r^{-1}\sum_{i=0}^{L_r}\delta_{t_{r,i}}\).

\begin{theorem}[Límite de la publicación nodal]
\label{thm:dim-nodal-limit}
Se cumple \(\nu_r\rightharpoonup\mu_K\). Si \(f\) es continua,
\[
 \left|\int f\,d\nu_r-\int f\,d\mu_K\right|
 \leq \omega_f(h_r)+\frac{\operatorname{osc}(f)}{N_r},
\]
donde \(\omega_f\) es su módulo de continuidad. Para \(p\geq1\),
el error del momento de orden \(p\) está acotado por
\(p h_r+1/N_r\). Para cada par fijo \((k,m)\),
\[
 \frac1{N_r}Y_r\longrightarrow (M_{a+b})_{
     0\leq a<k,\,0\leq b<k+m}=:Y_\infty^{(k,m)},
 \qquad\operatorname{rank}Y_\infty^{(k,m)}=k.
\]
\end{theorem}
\begin{proof}
La medida que asigna masa \(1/L_r\) a cada extremo izquierdo de
celda aproxima su integral con error a lo sumo \(\omega_f(h_r)\).
En efecto, cada celda tiene esa masa y diámetro no mayor que \(h_r\).
La medida \(\nu_r\) es \(L_r/N_r\) veces dicha medida más
\(\delta_1/N_r\), lo que añade a lo sumo
\(\operatorname{osc}(f)/N_r\). Para \(f(t)=t^p\), la constante
de Lipschitz es \(p\). La convergencia de cada entrada se sigue de
estas cotas. El bloque inicial de la matriz límite es definido positivo:
la integral del cuadrado de un polinomio no nulo es positiva, porque
\(\mu_K\) tiene densidad estrictamente positiva en todo intervalo
abierto de \([0,1]\). Así el límite conserva rango \(k\).
\end{proof}

Los momentos nodales varían exactamente por la incorporación de nodos.
Si \(\mathcal J_r\) es el conjunto de las nuevas posiciones y
\(z_q(t)=(1,t,\ldots,t^{q-1})^{\mathsf T}\), entonces
\[
 V_{r+1}^{(q)}(V_{r+1}^{(q)})^{\mathsf T}
 =V_r^{(q)}(V_r^{(q)})^{\mathsf T}
  +\sum_{i\in\mathcal J_r}z_q(t_{r+1,i})z_q(t_{r+1,i})^{\mathsf T}.
\]
Cada sumando es positivo semidefinido. Para las matrices rectangulares
\(Y_r\) vale la misma actualización con
\(z_kz_{k+m}^{\mathsf T}\). La igualdad expresa el cambio de la
medida de conteo; la normalización y el teorema anterior describen su
límite. El objeto límite de rango \(k\) tiene coordenadas en un
Grassmanniano finito, mientras los datos externos \(Z_r\) pertenecen
a las distintas publicaciones de cada nivel.

\subsection{Promedios celulares, operadores compatibles y detalle ortogonal}

Para una celda \(I_{r,i}=[b_{r,i},b_{r,i}+\ell_{r,i}]\), ponemos
\[
 a_{r,i}^{(p)}=
 \frac{(b_{r,i}+\ell_{r,i})^{p+1}-b_{r,i}^{p+1}}
 {(p+1)\ell_{r,i}},\qquad
 A_r^{(q)}=(a_{r,i}^{(p)})_{0\leq i<L_r,\,0\leq p<q}.
\]
Son las medias de los monomios respecto de la probabilidad condicional
de \(\mu_K\) en cada celda. Esta probabilidad es uniforme en la
celda, aunque \(\mu_K\) no sea la medida de longitud global.

Sea \(\mathcal H_r=\mathbb R^{L_r}\) con producto
\(\langle x,y\rangle_r=L_r^{-1}\sum_i x_i y_i\). La aplicación
\(J_r:\mathcal H_r\to\mathcal H_{r+1}\) repite el valor de un
padre en sus nueve hijos. Su adjunta \(R_r=J_r^*\) promedia esos
nueve valores. Se cumple \(R_rJ_r=I\).

\begin{theorem}[Compatibilidad exacta y balance de los momentos]
\label{thm:dim-cell-balance}
Para todo \(p\geq0\),
\[
 R_r a_{r+1}^{(p)}=a_r^{(p)},\qquad
 \frac1{L_r}\sum_i a_{r,i}^{(p)}=M_p.
\]
Definamos
\[
 D_r^{(q)}=A_{r+1}^{(q)}-J_rA_r^{(q)},\quad
 G_r^{(q)}=\frac1{L_r}(A_r^{(q)})^{\mathsf T}A_r^{(q)}.
\]
Entonces
\[
 R_rD_r^{(q)}=0,\qquad
 G_{r+1}^{(q)}=G_r^{(q)}+
       \frac1{L_{r+1}}(D_r^{(q)})^{\mathsf T}D_r^{(q)},
\]
y
\[
 G_r^{(q)}\longrightarrow H_K^{(q)}=(M_{a+b})_{0\leq a,b<q}.
\]
Los detalles y el nivel inicial reconstruyen todos los niveles. Para
\(q\leq L_r\), la matriz \(A_r^{(q)}\) tiene rango \(q\) y
todos los menores ordenados de orden \(q\) son positivos.
\end{theorem}
\begin{proof}
Las masas de los nueve hijos son iguales y sus intervalos particionan
el intervalo padre. La aditividad de la integral da la primera identidad;
la suma sobre todas las celdas da la segunda. La descomposición
\(A_{r+1}=J_rA_r+D_r\) es ortogonal en cada columna porque
\(R_rD_r=0\). Al expandir su Gram se anulan los términos cruzados,
y \(J_r\) es una isometría. Esto prueba el balance matricial.

El observable constante por celdas que toma el valor \(a_{r,i}^{(p)}\)
aproxima uniformemente \(t^p\), con error no mayor que \(p h_r\)
para \(p\geq1\). Por tanto converge en \(L^2(\mu_K)\), y los
productos internos convergen a \(M_{a+b}\). La reconstrucción se
obtiene iterando \(A_{r+1}=J_rA_r+D_r\).

Para probar los menores, elijamos \(q\) celdas ordenadas. Por
multilinealidad y Fubini, el determinante de sus promedios es la media
sobre el producto de esas celdas de
\(\prod_{a<b}(x_b-x_a)\). Este integrando es positivo en el
interior del producto porque las celdas son disjuntas y están ordenadas.
Sus fronteras tienen medida nula. La media es estrictamente positiva.
\end{proof}

La identidad se prolonga al par de grados \(k\) y \(k+m\):
\[
 \overline Y_r=\frac1{L_r}(A_r^{(k)})^{\mathsf T}A_r^{(k+m)},
 \qquad
 \overline Y_{r+1}-\overline Y_r
 =\frac1{L_{r+1}}(D_r^{(k)})^{\mathsf T}D_r^{(k+m)}.
\]
Para \(k+m\leq L_r\), las matrices de promedios definen una
segunda realización amplituhedral positiva, de \(L_r\) columnas,
y \(\overline Y_r\) tiene rango \(k\). Son promedios de celdas,
mientras \(Y_r\) usa evaluaciones en \(N_r=L_r+1\) nodos. Ambas
realizaciones convergen a \(Y_\infty^{(k,m)}\); su coincidencia
en el límite procede de las demostraciones anteriores.

La versión operatoria es igualmente explícita. Para una función continua
\(f\), sea \(\Psi_r(f)\) el operador diagonal de sus medias
condicionales en \(\mathcal H_r\). Es lineal, positivo y unital, y
\[
 R_r\Psi_{r+1}(f)J_r=\Psi_r(f).
\]
Esta igualdad se prueba aplicando ambos miembros a cada vector base.
Para \(X(t)=t\), los dos primeros operadores determinan, celda a
celda,
\[
 \Psi_r(X)_i=b_{r,i}+\frac{\ell_{r,i}}2,\qquad
 \Psi_r(X^2)_i-\Psi_r(X)_i^2=\frac{\ell_{r,i}^2}{12}.
\]
Si \(v_{r,i}\) denota esa varianza, la fórmula
\(\ell_{r,i}=\sqrt{12v_{r,i}}\) recupera la longitud; la
primera identidad recupera el extremo inicial. Con orientación, marcas y \(Q\),
se recupera nuevamente \(K\). El promedio es una compresión positiva:
la varianza escrita cuantifica por qué \(\Psi_r(X^2)\) y
\(\Psi_r(X)^2\) son operaciones distintas.

\subsection{Transporte de la medida y alcance del crecimiento dimensional}

Dos listas positivas de separaciones \(d\) y \(d'\), con las
mismas marcas, determinan un homeomorfismo creciente afín por intervalos
\(F=F_{K'}\circ F_K^{-1}\). Lleva cada celda a su celda homóloga
y satisface \(F_*\mu_K=\mu_{K'}\). Por ello
\[
 \mathcal U_F:L^2(\mu_K)\longrightarrow L^2(\mu_{K'}),
 \qquad\mathcal U_F f=f\circ F^{-1}
\]
es una isometría sobreyectiva. Si \(P_r\) y \(P'_r\) son las
proyecciones a funciones constantes por celdas, entonces
\(P'_r\mathcal U_F=\mathcal U_FP_r\): las medias condicionales
se transportan porque las masas de las celdas correspondientes coinciden.
Esto aplica en particular al flujo positivo de separaciones dirigido por
\(K\), conservando la dirección de cada bloque en sus hijos. Los
momentos de la coordenada espacial pueden cambiar; la isometría transporta
el observable junto con la medida y sus proyecciones.

El crecimiento demostrado es cofinal en toda dimensión externa finita.
Conserva un sistema de banderas positivas, una familia explícita de
puntos amplituhedrales, cuerpos convexos completos en rango uno y una
medida cuyos operadores de promedio son exactamente compatibles. La
memoria de los detalles permite invertir la compresión entre niveles.
Las pruebas de cobertura y residuos de rango uno conservan su dominio;
para rangos mayores, el resultado presente fija matrices, positividad,
recuperación, compatibilidad y límite con tipos explícitos. Esta
separación sitúa el alcance de cada construcción en el mismo desarrollo
sin confundir crecimiento de dimensión, parametrización de un punto y
clasificación de todas las celdas de una imagen.
